% GENERATED FILE. Edit canonical lessons, then run npm run generate:books.
% canonical-source-hash: fnv1a64:b140eac136538767
% canonical-lessons: KA-C423-hanta, KA-C423-karya, KA-C423-tantra, KA-C423-upaya, KA-C423-vibhajisu

\chapter{Step by Step}
\label{ch:ka-step-by-step}
\hlchaptermodality{423}

\begin{chapteropening}
\textbf{By the end of this chapter:} \emph{I can say a stage, a task, a strategy, a trick and to divide.}

Complete the last lesson of chapter 423: I can say a stage, a task, a strategy, a trick and to divide.
\end{chapteropening}

\section[\texorpdfstring{\kn{ಹಂತ} (hanta)}{ಹಂತ (hanta)}]{\kn{ಹಂತ} (hanta) — a stage; a phase}

\label{lesson:KA-C423-hanta}



\begin{quote}
Before the new one: say the Kannada for a method, then the Kannada for a process.
\end{quote}

\subsection*{You'll want to know: \kn{ಹಂತ}}
\textbf{\kn{ಹಂತ}} — \emph{hanta} — \textquotedblleft{}a stage; a phase\textquotedblright{}.

\begin{quote}
\kn{ಕೆಲಸದ} \kn{ಮೊದಲ} \kn{ಹಂತ} \kn{ಮುಗಿಯಿತು}.
\end{quote}

\begin{quote}
\emph{kelasada modala hanta mugiyitu.}
\end{quote}

\begin{quote}
The first stage of the work is over.
\end{quote}

\kn{ಹಂತಹಂತವಾಗಿ}, step by step.

\begin{cousinweb}
\textbf{\kn{ಹಂತ}} (\emph{hanta}) is a Kannada word whose origin is not settled: neither DEDR nor CDIAL lists it.
\end{cousinweb}

\begin{grammarlens}[title={What you've built}]
One more word for saying what happens if.
\end{grammarlens}

\subsection*{Your turn}
\begin{itemize}
\raggedright
  \item \emph{Say it:} \emph{hanta}
  \item \emph{Say it:} \emph{hanta}, once more
  \item \emph{Say it:} \emph{kelasada modala hanta mugiyitu.} — The first stage of the work is over.
  \item \emph{Recall:} say the Kannada for a method, then the Kannada for a process, then say \emph{hanta} again
\end{itemize}

\begin{tcolorbox}[breakable,colback=teal!4,colframe=teal!35!black,title={Before you move on}]
What does \kn{ಹಂತ} mean? (\textquotedblleft{}A stage; a phase\textquotedblright{}.) Say it once more.
\end{tcolorbox}
\section[\texorpdfstring{\kn{ಕಾರ್ಯ} (kārya)}{ಕಾರ್ಯ (kārya)}]{\kn{ಕಾರ್ಯ} (kārya) — a task; an act}

\label{lesson:KA-C423-karya}



\begin{quote}
Before the new one: say the Kannada for a process, then the Kannada for a stage.
\end{quote}

\subsection*{You'll want to know: \kn{ಕಾರ್ಯ}}
\textbf{\kn{ಕಾರ್ಯ}} — \emph{kārya} — \textquotedblleft{}a task; an act\textquotedblright{}.

\begin{quote}
\kn{ಈ} \kn{ಕಾರ್ಯಕ್ಕೆ} \kn{ಎಲ್ಲರ} \kn{ಸಹಾಯ} \kn{ಬೇಕು}.
\end{quote}

\begin{quote}
\emph{ī kāryakke ellara sahāya bēku.}
\end{quote}

\begin{quote}
This task needs everyone's help.
\end{quote}

\kn{ಕಾರ್ಯಕ್ರಮ}, a programme, is built on it.

\begin{cousinweb}
\textbf{\kn{ಕಾರ್ಯ}} (\emph{kārya}) is borrowed from Sanskrit \emph{kārya}, \textquotedblleft{}to be done; a task, a duty\textquotedblright{} (CDIAL 3078).
\end{cousinweb}

\begin{grammarlens}[title={What you've built}]
One more word for saying what happens if.
\end{grammarlens}

\subsection*{Your turn}
\begin{itemize}
\raggedright
  \item \emph{Say it:} \emph{kārya}
  \item \emph{Say it:} \emph{kārya}, once more
  \item \emph{Say it:} \emph{ī kāryakke ellara sahāya bēku.} — This task needs everyone's help.
  \item \emph{Recall:} say the Kannada for a process, then the Kannada for a stage, then say \emph{kārya} again
\end{itemize}

\begin{tcolorbox}[breakable,colback=teal!4,colframe=teal!35!black,title={Before you move on}]
What does \kn{ಕಾರ್ಯ} mean? (\textquotedblleft{}A task; an act\textquotedblright{}.) Say it once more.
\end{tcolorbox}
\section[\texorpdfstring{\kn{ತಂತ್ರ} (tantra)}{ತಂತ್ರ (tantra)}]{\kn{ತಂತ್ರ} (tantra) — a strategy; a technique}

\label{lesson:KA-C423-tantra}



\begin{quote}
Before the new one: say the Kannada for a stage, then the Kannada for a task.
\end{quote}

\subsection*{You'll want to know: \kn{ತಂತ್ರ}}
\textbf{\kn{ತಂತ್ರ}} — \emph{tantra} — \textquotedblleft{}a strategy; a technique\textquotedblright{}.

\begin{quote}
\kn{ಆಟ} \kn{ಗೆಲ್ಲಲು} \kn{ನಮ್ಮ} \kn{ತಂಡ} \kn{ಹೊಸ} \kn{ತಂತ್ರ} \kn{ಬಳಸಿತು}.
\end{quote}

\begin{quote}
\emph{āṭa gellalu namma taṇḍa hosa tantra baḷasitu.}
\end{quote}

\begin{quote}
Our team used a new strategy to win the game.
\end{quote}

\kn{ತಂತ್ರಜ್ಞಾನ}, technology.

\begin{cousinweb}
\textbf{\kn{ತಂತ್ರ}} (\emph{tantra}) is borrowed from Sanskrit \emph{tantra}, \textquotedblleft{}a loom, the warp; a system, a device\textquotedblright{} (CDIAL 5663).
\end{cousinweb}

\begin{grammarlens}[title={What you've built}]
One more word for saying what happens if.
\end{grammarlens}

\subsection*{Your turn}
\begin{itemize}
\raggedright
  \item \emph{Say it:} \emph{tantra}
  \item \emph{Say it:} \emph{tantra}, once more
  \item \emph{Say it:} \emph{āṭa gellalu namma taṇḍa hosa tantra baḷasitu.} — Our team used a new strategy to win the game.
  \item \emph{Recall:} say the Kannada for a stage, then the Kannada for a task, then say \emph{tantra} again
\end{itemize}

\begin{tcolorbox}[breakable,colback=teal!4,colframe=teal!35!black,title={Before you move on}]
What does \kn{ತಂತ್ರ} mean? (\textquotedblleft{}A strategy; a technique\textquotedblright{}.) Say it once more.
\end{tcolorbox}
\section[\texorpdfstring{\kn{ಉಪಾಯ} (upāya)}{ಉಪಾಯ (upāya)}]{\kn{ಉಪಾಯ} (upāya) — a trick; a way out}

\label{lesson:KA-C423-upaya}



\begin{quote}
Before the new one: say the Kannada for a task, then the Kannada for a strategy.
\end{quote}

\subsection*{You'll want to know: \kn{ಉಪಾಯ}}
\textbf{\kn{ಉಪಾಯ}} — \emph{upāya} — \textquotedblleft{}a trick; a way out\textquotedblright{}.

\begin{quote}
\kn{ಮಳೆಯಿಂದ} \kn{ತಪ್ಪಿಸಿಕೊಳ್ಳಲು} \kn{ಅವನಿಗೆ} \kn{ಒಂದು} \kn{ಉಪಾಯ} \kn{ಹೊಳೆಯಿತು}.
\end{quote}

\begin{quote}
\emph{maḷeyinda tappisikoḷḷalu avanige ondu upāya hoḷeyitu.}
\end{quote}

\begin{quote}
A way to escape the rain came to his mind.
\end{quote}

\kn{ಉಪಾಯ} \kn{ಹೊಳೆ}, for an idea to strike.

\begin{cousinweb}
\textbf{\kn{ಉಪಾಯ}} (\emph{upāya}) is borrowed from Sanskrit \emph{upāya}, \textquotedblleft{}coming near; a means, a way\textquotedblright{} (CDIAL 2308).
\end{cousinweb}

\begin{grammarlens}[title={What you've built}]
One more word for saying what happens if.
\end{grammarlens}

\subsection*{Your turn}
\begin{itemize}
\raggedright
  \item \emph{Say it:} \emph{upāya}
  \item \emph{Say it:} \emph{upāya}, once more
  \item \emph{Say it:} \emph{maḷeyinda tappisikoḷḷalu avanige ondu upāya hoḷeyitu.} — A way to escape the rain came to his mind.
  \item \emph{Recall:} say the Kannada for a task, then the Kannada for a strategy, then say \emph{upāya} again
\end{itemize}

\begin{tcolorbox}[breakable,colback=teal!4,colframe=teal!35!black,title={Before you move on}]
What does \kn{ಉಪಾಯ} mean? (\textquotedblleft{}A trick; a way out\textquotedblright{}.) Say it once more.
\end{tcolorbox}
\section[\texorpdfstring{\kn{ವಿಭಜಿಸು} (vibhajisu)}{ವಿಭಜಿಸು (vibhajisu)}]{\kn{ವಿಭಜಿಸು} (vibhajisu) — to divide; to split}

\label{lesson:KA-C423-vibhajisu}



\begin{quote}
Before the new one: say the Kannada for a strategy, then the Kannada for a trick.
\end{quote}

\subsection*{You'll want to know: \kn{ವಿಭಜಿಸು}}
\textbf{\kn{ವಿಭಜಿಸು}} — \emph{vibhajisu} — \textquotedblleft{}to divide; to split\textquotedblright{}.

\begin{quote}
\kn{ಶಿಕ್ಷಕರು} \kn{ತರಗತಿಯನ್ನು} \kn{ನಾಲ್ಕು} \kn{ಗುಂಪುಗಳಾಗಿ} \kn{ವಿಭಜಿಸಿದರು}.
\end{quote}

\begin{quote}
\emph{śikṣakaru taragatiyannu nālku gumpugaḷāgi vibhajisidaru.}
\end{quote}

\begin{quote}
The teacher divided the class into four groups.
\end{quote}

Past (I …): \kn{ವಿಭಜಿಸಿದೆ}, \emph{vibhajiside}.

Into what: -\kn{ಆಗಿ}, as in \kn{ಗುಂಪುಗಳಾಗಿ}.

\begin{cousinweb}
\textbf{\kn{ವಿಭಜಿಸು}} (\emph{vibhajisu}) is built on Sanskrit \emph{vibhajana}, \textquotedblleft{}separation, distinction\textquotedblright{} (Monier-Williams), with the Kannada verb ending -\emph{isu}.
\end{cousinweb}

\begin{grammarlens}[title={What you've built}]
One more word for saying what happens if.
\end{grammarlens}

\subsection*{Your turn}
\begin{itemize}
\raggedright
  \item \emph{Say it:} \emph{vibhajisu}
  \item \emph{Say it:} \emph{vibhajisu}, once more
  \item \emph{Say it:} \emph{śikṣakaru taragatiyannu nālku gumpugaḷāgi vibhajisidaru.} — The teacher divided the class into four groups.
  \item \emph{Recall:} say the Kannada for a strategy, then the Kannada for a trick, then say \emph{vibhajisu} again
\end{itemize}

\begin{tcolorbox}[breakable,colback=teal!4,colframe=teal!35!black,title={Before you move on}]
What does \kn{ವಿಭಜಿಸು} mean? (\textquotedblleft{}To divide; to split\textquotedblright{}.) Say it once more.
\end{tcolorbox}
